% The BMJ: verified 2022 Research layout, neutral Chinese adaptation.
\PassOptionsToPackage{no-math}{fontspec}
\documentclass[UTF8,fontset=none,scheme=plain]{ctexart}
\usepackage[paperwidth=595.276bp,paperheight=793.701bp,left=34bp,right=34bp,top=64bp,bottom=40bp,headheight=17bp,headsep=15bp,footskip=18bp]{geometry}
% Shared dependencies and neutral material. Compile with XeLaTeX.
% No publisher logos, paper text, data, figures, or download metadata are included.
\usepackage{amsmath,amssymb,graphicx,xcolor,tikz,fancyhdr,needspace,xurl,hyperref,balance,etoolbox}
% The legacy balance package omits the first split box's final baseline depth.
% Reserve that measured depth rather than hiding a CJK overfull-box warning.
\makeatletter
\patchcmd{\@BAlancecol}{\setbox\@leftcolumn=\vbox to \@tempdima}{\advance\@tempdima by \dp1\relax\setbox\@leftcolumn=\vbox to \@tempdima}{}{\PackageWarning{cns-template}{Balance depth patch unavailable; inspect final columns}}
\makeatother
\usetikzlibrary{arrows.meta,positioning}
\usepackage[protrusion=true,expansion=false]{microtype}
\IfFontExistsTF{Noto Serif CJK SC}{%
  \setCJKmainfont{Noto Serif CJK SC}%
}{%
  \setCJKmainfont[BoldFont=FandolSong-Bold.otf,ItalicFont=FandolKai-Regular.otf]{FandolSong-Regular.otf}%
}
\IfFontExistsTF{Noto Sans CJK SC}{%
  \setCJKsansfont{Noto Sans CJK SC}%
}{%
  \setCJKsansfont[BoldFont=FandolHei-Bold.otf]{FandolHei-Regular.otf}%
}
\setmainfont{TeX Gyre Termes}\setsansfont{TeX Gyre Heros}
\setmonofont{TeX Gyre Cursor}
\definecolor{JournalBlue}{HTML}{0097C9}
\definecolor{JournalRed}{HTML}{B83443}
\definecolor{ScienceRed}{HTML}{CD1021}
\definecolor{CiteBlue}{HTML}{0000EE}
\hypersetup{colorlinks=true,linkcolor=JournalBlue,urlcolor=JournalBlue,
  pdftitle={CNS论文汉化：中性版式示例},pdfauthor={示例模板；非期刊官方模板},
  pdfsubject={Neutral source-inspired Chinese layout demonstration; no real research claims}}
\urlstyle{same}\raggedbottom
\setlength{\parindent}{1em}\setlength{\parskip}{0pt}
\setlength{\emergencystretch}{2em}\tolerance=2200
\XeTeXlinebreaklocale "zh"\XeTeXlinebreakskip=0pt plus .1pt
\pagestyle{fancy}\fancyhf{}
\renewcommand{\headrulewidth}{0pt}\renewcommand{\footrulewidth}{0pt}
\newcommand{\PaperTitle}{中文题名占位：从组织结构到证据解释}
\newcommand{\SampleAuthors}{Author A\textsuperscript{1}, Author B\textsuperscript{2,*}}
\newcommand{\SampleAffiliations}{\textsuperscript{1}作者单位一（示例占位）。\quad
  \textsuperscript{2}作者单位二（示例占位）。\par
  *通信作者信息待填。DOI、收稿和出版日期均应从目标原文核实。}
\newcommand{\SampleAbstract}{本段是中性摘要占位文字，不陈述任何真实研究发现。它演示中文题名、作者信息、摘要和双栏正文之间的层级关系。实际使用时，应完整替换为目标论文的忠实中文译文，并保留研究对象、方法、结果和结论的限定条件。图注、参考文献、作者声明和来源信息均属于独立的覆盖项。模板中的图形仅说明排版结构，不表示生物学机制、实验结果或统计关系。}
\newcommand{\DemoCaption}{\textbf{图1. 中性证据链示意（占位图，无实验数据）。}
  \textbf{A}，输入节点表示需要核对的原文材料。
  \textbf{B}，处理节点表示逐段转换及复核。
  \textbf{C}，输出节点表示正文、图注和书目等独立交付部分。
  箭头仅表示本示例的阅读顺序，不代表因果关系、效应方向或数量变化。
  该图由模板自行绘制，所有节点均为文字占位；无样本量、误差线、统计检验或比例尺。
  实际论文图注应逐项保留原稿报告的样本量、单位、检验方法、限定及图形署名。}
\newcommand{\DemoVector}{%
 \resizebox{\linewidth}{!}{%
 \begin{tikzpicture}[font=\sffamily\small,node distance=18mm,
   box/.style={draw=black!70,rounded corners=2pt,fill=black!3,minimum width=26mm,minimum height=16mm,align=center},
   flow/.style={-{Stealth[length=2.5mm]},line width=.7pt}]
  \node[box] (a) {\textbf{A}\quad 输入\\原文材料};
  \node[box,right=of a] (b) {\textbf{B}\quad 处理\\转换与复核};
  \node[box,right=of b] (c) {\textbf{C}\quad 输出\\中文译稿};
  \draw[flow] (a)--(b);\draw[flow] (b)--(c);
  \node[below=11mm of a,align=center] {标题与摘要\\正文与图表};
  \node[below=11mm of b,align=center] {保留限定\\检查单位};
  \node[below=11mm of c,align=center] {可追溯引文\\明确原文信息};
  \node[below=29mm of b,font=\sffamily\footnotesize,text=black!65] {示意占位；不包含真实实验数据};
 \end{tikzpicture}}%
}
\newcommand{\SampleBibliography}{%
 \section*{参考文献（示例占位）}
 {\fontsize{7.6bp}{10bp}\selectfont\setlength{\parindent}{0pt}\interlinepenalty=10000
 \hypertarget{ref:1}{}1. Author A. (20XX). Example reference title. \textit{Journal placeholder}. 卷、页码与DOI待填。此条仅供演示引文跳转。\par\smallskip
 \hypertarget{ref:2}{}2. Author B. (20YY). Another example reference title. \textit{Journal placeholder}. 出版信息待填；不要把占位条目带入正式译稿。\par}
}
\newcommand{\SampleSourceInfo}{%
 \section*{原文信息}
 以下字段全部为示例占位：作者；英文原标题；期刊；年份；卷期；页码或文章号；DOI或来源链接。实际使用时必须逐项替换并核对。\par
 本文件是\JournalName 风格的非官方中文译稿样板，依据一次已核验文章版式进行结构化改编，不代表该期刊所有年代或文体的统一标准，也不保证像素级复刻。无真实论文内容或期刊标识图像。\par
}
% Deliberately neutral demo prose, not translated from a published paper.
\newcommand{\SampleOpening}{%
 \section*{引言：示例正文与引文}
 这是一段用于检查双栏阅读节奏的中性示例文字。段落中的术语、缩写和标点需要在实际项目中与源文保持一致。中文译稿应呈现原论文的完整论述，而不是把多段证据压缩成一段摘要。这里的引文只链接到本文件末尾明确标为占位的参考文献\SampleCiteOne。\par
 若原句仅表示相关、可能或提示，译文也应保留相同的证据强度。数字、单位、正负号、上下标和范围端点应逐项核对。原文存在内部矛盾时，应分别保留相应位置的陈述，再用简短译注说明；不能根据常识把不同的数字静默统一。\par
 \subsection*{内容层级与跨页衔接}
 此小标题演示二级层级。实际论文可能采用不同的小标题密度，模板不会替代对原文结构的识别。双栏文章需按左栏、右栏和跨页续句的顺序读取。若段落在页末中断，应在完成整句核对后接入下一页，避免把图中文字误当正文，或重复翻译同一段。\par
 图表具有独立的信息范围。正文引用到某幅图，并不意味着其图注可以省略。图注中的每个分面、比较对象、时间点、样本量及统计说明都应对应源文。英文图内标签可以保留，并在图旁或文末提供中文释义；实际采用哪种安排，需要依据图形尺寸和译文长度判断\SampleCiteTwo。\par
 \subsection*{来源与可维护数据}
 本段继续演示正文密度。可编辑文本和图注数据应与最终版面分开维护，以便修订时准确定位变化。若只调整版式，应能证明正文、图注和书目内容没有被意外改写。更换字体或纸张后，须重新检查断行、图注与图片的相对位置以及页脚可读性。\par
 文内引用与书目之间应有明确对应关系。这里的两条演示书目没有真实DOI，也没有虚构的实验信息。使用模板时可增删条目，但必须同时更新编号、链接目标和引用位置。页码只表示中文译稿自己的分页，不用于冒充原出版页码。\par
 \subsection*{结果叙述的演示占位}
 本节的文字仅用于呈现长段落在双栏页面中的流动，不报告真实结果。实际翻译应保留观察对象、实验条件和结局指标之间的关系。例如，原文分别描述体外实验与动物实验时，不能把两类结果合并为同一层面的结论；研究使用的对照、时间顺序和分析范围也必须保留。图表中的比较范围不应扩大为对整个人群或全部疾病阶段的判断。\par
 论述中的转折与限制同样属于正文。若研究先提出一种解释，再讨论另一种可能，译文应沿用这一推进顺序。作者明确写出的不确定性应能在中文中被识别；如果尚未确定某个现象的方向或机制，就不能改成已经建立的因果关系。对于术语近似但研究层级不同的概念，可在首次出现时保留缩写，后文保持统一。\par
 \subsection*{方法信息与可复核性}
 方法部分往往包含密集数字和专有名词。翻译时须区分浓度、总量、频率和处理时长，不能因为中文句式调整而改变参数的归属。试剂名称、软件版本和数据访问标识通常保留原样；上下限、排除规则和显著性标准则需要逐字符核对。这里没有提供任何真实实验参数，避免把模板占位误认为可以直接执行的方法。\par
 表格与补充材料的引用也需要保留原层级。正文对补充图的指向并不意味着模板中已经附有该补充图；若源文件只包含主文，应如实标注所交付的范围。图表编号、附录顺序和参考文献数量都应由目标文章决定，本示例中的结构不构成固定配额，也不要求所有论文采用同样的章节。\par
}
\newcommand{\SampleClosing}{%
 \section*{讨论：版式适配与范围}
 上方矢量图是模板原创的排版占位图，不来自任何论文。图形与完整演示图注被放在同一页，便于读者直接核对分面说明。真实项目应使用经过授权或用户提供的原图素材，并保留图注署名与许可信息。尺寸不足时，可适度缩放或调整栏宽，不以删除图注解决空间问题。\par
 模板仅记录一次已核验文章的可观察设计特征。即使文章出自同一期刊，不同年代、研究论文与综述、校正清样与正式版本也可能采用不同的页眉和引文形式。应以用户提供的实际PDF为准，不能仅凭期刊名称推断配色、字号或引用方式\SampleCiteOne。\par
 \subsection*{演示检查项}
 本样板用于展示可编译性、中文字形、标题层级、跨栏图注、页眉页脚和内部链接。图中没有实际数据，因此不应被当作科学插图引用。替换内容后，还需检查所有页的渲染结果，确认无截字、碰撞、异常空白和遗漏段落。\par
 \SampleBibliography
 \SampleSourceInfo
}


\usepackage{multicol,changepage,booktabs,array}
\newcommand{\JournalName}{The BMJ}
\definecolor{BMJBlue}{HTML}{0070B9}
\hypersetup{linkcolor=black,urlcolor=BMJBlue}
\setlength{\columnsep}{13bp}
\renewcommand{\normalsize}{\fontsize{8.8bp}{11.6bp}\selectfont}
\ctexset{section={format=\sffamily\bfseries\fontsize{10.5bp}{13bp}\selectfont,beforeskip=8bp,afterskip=3bp},subsection={format=\sffamily\bfseries,beforeskip=6bp,afterskip=2bp}}
\newcommand{\SampleCiteOne}{\textsuperscript{\hyperlink{ref:1}{1}}}
\newcommand{\SampleCiteTwo}{\textsuperscript{\hyperlink{ref:2}{2}}}
\fancyhead[C]{\colorbox{BMJBlue}{\makebox[\dimexpr\textwidth-2\fboxsep\relax][r]{\color{white}\sffamily\bfseries\fontsize{9bp}{11bp}\selectfont RESEARCH}}}
\fancyfoot[L]{\fontsize{6.5bp}{8bp}\selectfont\textcolor{BMJBlue}{\textbf{the bmj}}\,|\,年份、卷及文章号待填\,|\,非官方中性样板}
\fancyfoot[R]{\sffamily\fontsize{6.8bp}{8bp}\selectfont 中文译页\quad\textbf{\thepage}}
\newcommand{\BMJLabel}[1]{\par\smallskip\noindent\textcolor{BMJBlue}{\sffamily\bfseries #1}\par\nobreak\noindent}
\newenvironment{BMJColumns}{\begin{adjustwidth}{107bp}{0bp}\begin{multicols}{2}\normalsize}{\end{multicols}\end{adjustwidth}}
\begin{document}\normalsize
\begin{tikzpicture}[remember picture,overlay]
\node[anchor=north west,inner sep=0,text width=96bp] at ([xshift=34bp,yshift=-68bp]current page.north west){\sffamily\fontsize{7bp}{9bp}\selectfont
\textbf{OPEN ACCESS}\par 开放获取占位\par\vspace{65bp}
编号作者单位见文末。\par\smallskip 通信作者与联系方式待填。\par\smallskip
\textcolor{BMJBlue}{引用、DOI及接受日期待核实。}\par\smallskip 本栏演示原刊首页的外侧元数据窄栏，不包含真实作者或出版记录。};
\end{tikzpicture}
\begin{adjustwidth}{107bp}{0bp}
\noindent\rule{\linewidth}{.35bp}\par\vspace{3bp}
{\sffamily\fontsize{17bp}{20.5bp}\selectfont\noindent 从原文材料到可追溯译稿：结构化证据与版式核对的中性演示\par}\vspace{7bp}
{\sffamily\fontsize{8.5bp}{11.3bp}\selectfont\noindent\SampleAuthors\quad（作者占位）\par}\vspace{10bp}
\noindent\begin{minipage}[t]{203bp}\sffamily\fontsize{8.1bp}{10.5bp}\selectfont\setlength{\parindent}{0pt}
\textbf{摘要}
\BMJLabel{目的}演示如何保留研究问题、比较方向和来源范围。本文件不含临床研究结论。
\BMJLabel{设计}中性结构化版式占位，不对应真实试验。
\BMJLabel{数据来源}实际翻译应以用户提供的原始文件为依据；本样板没有真实数据来源或检索结果。
\BMJLabel{研究纳入标准}此处演示资格标准标签。目标文章的筛选条件、样本和时间范围必须完整翻译，不能从模板生成。
\BMJLabel{主要结局指标}保留原文关于结局、统计方法和不确定性的定义。若原表用符号表达证据等级，中文版需要给出清楚的释义。本文只演示段落密度，不含效应量、置信区间或确定性评价。\par
摘要中的研究对象、干预和比较应与源文逐项对应。实际项目需要区分相对效应与绝对效应、单药与联合用药、基线风险与观察到的事件数。仅在原文提供依据时，才记录相应统计方法和比较方向。\par\smallskip
对于需要换样式的旧译稿，应保存原译文的内容记录，再检查题名、摘要、正文、图表脚注和书目是否保持完整。排版参数可以调整，出版字段和科学内容则必须逐项对应当前来源版本。
\par\vspace{12bp}
\makebox[\linewidth][l]{\hspace*{-107bp}\fcolorbox{BMJBlue!50}{white}{\begin{minipage}{\dimexpr310bp-2\fboxsep-2\fboxrule\relax}\color{BMJBlue}\fontsize{7.7bp}{10bp}\selectfont
\textbf{本主题已有认识（占位）}\par 完整译稿包含摘要、正文、图表、脚注、声明和书目；这些内容需要分别登记并核对。\par\smallskip
\textbf{本示例新增内容（占位）}\par 蓝色框文向左伸入元数据侧栏，演示本次核验版本的首页结构。此处文字均为原创版式说明，不可作为科学证据引用。\par
图文应在相关讨论附近交替出现；图注变紧凑时，仍应保留所有面板和统计说明。\par\smallskip
在正式译稿中，本框应分别完整翻译已有认识与本研究新增认识。原刊使用的项目、限定和不确定性应逐项保留；示例框的段落数量不构成固定要求。
\end{minipage}}}
\end{minipage}\hfill\begin{minipage}[t]{203bp}\sffamily\fontsize{8.1bp}{10.5bp}\selectfont\setlength{\parindent}{0pt}
\BMJLabel{结果}\SampleAbstract\par\smallskip
本段仅展示较长结果摘要在右栏中的连续阅读。源文分别报告的比较和风险层级需要逐项对应，不能为了填满版面而虚构数值，也不能把不同人群的相对与绝对效应混为一谈。\par\smallskip
如原文存在不同位置的数值不一致，应保留各处记录并给出简短译注。可编辑表格必须能追溯到原表的行、列与脚注；表格截图不能替代需要汉化的中文表格。\par\smallskip
图内文字保留英文时，须给出完整中文释义。原刊版权和许可属于正常出版信息；个人下载来源则不混入正文或分发素材。\par\smallskip
长表在分页时重复表题和列头，使读者能直接辨认续页中的比较对象。原表共享的统计单元格如果在重排后重复展示，应在脚注中说明对应关系，避免将重复显示误读成新增分析。所有符号、等级和范围端点都应保留。\par\smallskip
双栏续段、跨页图注、附录指向和作者声明具有独立的阅读顺序。首次抽取后要结合页面图像检查，确认没有把页脚插入正文，也没有把右栏提前接在左栏未完成的句子中。译文长度改变分页时，原文的逻辑顺序仍应连贯。
\BMJLabel{结论}本页只展示一种已核验的首页结构。正式使用时应替换全部占位内容，并以当前原文核实出版年份、文章号与DOI。
\end{minipage}
\end{adjustwidth}
\clearpage
\begin{BMJColumns}\SampleOpening
\subsection*{长表与正文的阅读衔接}
表格宜靠近首次解释其主要结果的段落。若原表过宽，可以重组字段或分为连续部分，保留每一行的研究名称、单位、比较、随访和脚注。拆分后应保留明确的关联标识，并逐项核对，不以缩小到难以阅读的字体来维持原页数。\par
全宽图表会中断双栏正文。恢复正文时，应让读者能够识别前后续句；必要时调整段落断点或图像尺寸。结束页可平衡两栏，普通正文页则应检查是否存在因不合理强制分页造成的大块空白。\par
\subsection*{数字与版本字段的复核}
数字核对需要同时关注比较方向、正负号、时间、样本量、分母和计量单位。只比较数字字符的集合不能证明对应关系正确。原文将两项比较共用一个估计时，应记录这一结构，并在中文重排中清楚表示。\par
提前在线版本可能尚未出现最终卷期或页码。接受日期、在线日期、刊期和版权年份应分别记录；文件名中的年份不能直接替代原刊给出的出版信息。译文文末应独立列出完整原文信息，使读者能够追溯所使用的版本。
\end{BMJColumns}
\clearpage
\noindent\begin{minipage}{.57\linewidth}\DemoVector\end{minipage}\hfill\begin{minipage}{.39\linewidth}\sffamily\fontsize{7.6bp}{9.2bp}\selectfont\DemoCaption\end{minipage}\par\vspace{10bp}
{\sffamily\fontsize{8bp}{10.2bp}\selectfont\noindent\textbf{表1．中性可编辑证据记录（无真实数据）}\par\smallskip
\noindent\begin{tabular}{@{}p{110bp}p{170bp}p{220bp}@{}}\toprule
材料 & 核对字段 & 来源与范围\\\midrule
图表占位 & 标签、方向、单位、完整图注 & 原始文件中逐幅确认；不得继承旧裁区\\[4bp]
正文占位 & 跨页续句、否定、不确定性 & 恢复正确阅读顺序后逐段核对\\[4bp]
书目占位 & 编号、链接、英文书目信息 & 对应目标原文，非模板固定数量\\\bottomrule
\end{tabular}\par}
\vspace{8bp}
\begin{BMJColumns}
\section*{图、图注与标签释义}
本页的矢量图由模板原创，仅说明材料与处理之间的阅读关系。其尺寸没有实验意义，箭头没有因果含义。正式使用时需替换为目标原图，并逐幅核实所有面板、坐标轴和尺度，没有沿用某个旧项目的固定裁剪区域。\par
图注采用较小字号与行距，以保持图文紧凑。压缩前后应逐项比较面板说明、统计检验、样本量、误差线、缩写和署名；字体变小不能成为删去长注的理由。长标签对照可以接续在文末，并清楚对应图号。
\SampleClosing\end{BMJColumns}
\end{document}
